\section{FUNDAMENTAÇÃO TEÓRICA E TRABALHOS RELACIONADOS}

Este capítulo reúne os conceitos que orientaram as decisões de projeto e compara a Agendinha com serviços de agendamento disponíveis no mercado. A separação entre fundamentação e desenvolvimento permite que o capítulo seguinte se concentre exclusivamente no que foi implementado.

\subsection{Fundamentos técnicos}

\subsubsection{Digitalização enxuta e experiência mobile-first}

Digitalizar um processo não significa apenas converter o caderno em uma tela. O sistema deve retirar etapas manuais, centralizar o estado e tornar o resultado verificável. No caso estudado, o ganho esperado vem da publicação automática da disponibilidade e da validação da reserva, e não da inclusão de módulos administrativos. A abordagem mobile-first parte da menor largura de tela e da interação por toque, coerente com um público que pode ter o celular como principal dispositivo.

\subsubsection{Progressive Web Apps}

Uma Progressive Web App combina tecnologias da Web com recursos de instalação e execução semelhantes aos de aplicativos. Manifesto, HTTPS e service worker permitem adicionar a aplicação à tela inicial e controlar estratégias de rede e cache \cite{mdnpwa,w3cmanifest}. Essa alternativa evita manter projetos nativos separados e reduz a barreira de instalação. Não implica, porém, funcionamento integral sem internet: reservas precisam alcançar o servidor para que a concorrência seja resolvida sobre um estado comum.

\subsubsection{Arquitetura gerenciada e Cache-Aside}

Serviços gerenciados deslocam tarefas de provisionamento, atualização e disponibilidade para provedores especializados. Isso favorece equipes pequenas, embora introduza dependência externa. Na estratégia Cache-Aside, a aplicação consulta o Redis antes da fonte oficial; em caso de ausência, lê o PostgreSQL, grava uma cópia temporária e devolve o resultado \cite{fowler2002}. Uma versão por prestador invalida logicamente as chaves após mudanças. Se o Redis estiver indisponível, a consulta permanece correta, apenas potencialmente mais lenta.

\subsubsection{Transações, concorrência e idempotência}

O horário exibido ao cliente pode ser escolhido por outra pessoa antes da confirmação. Por isso, a verificação decisiva ocorre na transação de reserva. A restrição de exclusão do PostgreSQL impede dois intervalos confirmados sobrepostos para o mesmo prestador \cite{postgresconstraints}. A chave de idempotência associa tentativas repetidas à mesma resposta, evitando duplicação quando o cliente reenvia a requisição por perda de conexão.

\subsubsection{Segurança e comunicação assíncrona}

O Supabase Auth emite JWTs para identificar o prestador, enquanto as políticas de Row Level Security limitam linhas ao respectivo proprietário \cite{supabasetls}. O cliente não recebe acesso amplo às tabelas: usa funções públicas com validação e token específico para cancelamento. Para e-mails e avisos, o padrão Transactional Outbox grava o evento na mesma transação do domínio; workers realizam a entrega depois do commit e registram tentativas. Assim, uma falha do provedor de mensagens não desfaz uma reserva válida.

Esses conceitos conduzem a quatro propriedades observáveis: completude das tarefas, ausência de sobreposição, degradação segura sem cache e custo compatível com a demanda. É sobre essas propriedades, e não sobre a simples presença das tecnologias, que o projeto deve ser analisado.

\subsection{Trabalhos relacionados}

Foram selecionados três serviços de agendamento com páginas oficiais consultadas em 25 de outubro de 2026. A Tabela~\ref{tab:comparacao} contrapõe informações comerciais verificáveis ao escopo implementado na Agendinha. A presença de recursos adicionais é característica da oferta, não evidência de pior usabilidade. A comparação também não pressupõe que acesso por navegador signifique implementação como PWA.

\begin{table}[htbp]
\centering
\caption{Comparação de escopo e condições anunciadas por sistemas de agendamento.}
\label{tab:comparacao}
\footnotesize
\setlength{\tabcolsep}{4pt}
\begin{tabularx}{\textwidth}{YYYYY}
\toprule
\textbf{Critério} & \textbf{Agendinha} & \textbf{Booksy} & \textbf{Trinks} & \textbf{Reservio} \\
\midrule
\textbf{Oferta principal} & Agenda essencial para autônomos & Agenda e gestão para beleza e bem-estar & Agenda e gestão de negócios de beleza & Reservas e gestão por níveis de plano \\
\textbf{Preço divulgado} & Sem preço comercial definido & R\$ 99,99/mês + impostos & R\$ 110/mês (1--2 profissionais, mensal) & Plano Free: R\$ 0; planos pagos com preço variável \\
\textbf{Limite explícito} & Depende da capacidade testada & Agendamentos ilimitados anunciados & Não identificado na página consultada & Free: 40 reservas por 30 dias \\
\textbf{Agendamento web} & Página pública compartilhável & Agendamento online 24/7 & Site e app de agendamento & Site de reservas \\
\textbf{Recursos adicionais anunciados} & Não inclui pagamentos, estoque ou finanças & Marketing, pagamentos e relatórios & Gestão financeira, estoque e pagamentos & Estatísticas, calendário e recursos por plano \\
\textbf{Instalação PWA} & Prevista na aplicação & Não aferida & Não aferida & Não aferida \\
\bottomrule
\end{tabularx}
\vspace{2mm}

{\footnotesize\raggedright Fonte: elaboração própria, a partir das páginas comerciais oficiais \cite{booksypricing,trinkspricing,reserviopricing}. Consulta em 25 out. 2026. Os valores não incluem condições promocionais nem custos adicionais não informados.\par}
\end{table}

O Booksy anuncia R\$ 99,99 por mês, mais impostos, e cobrança adicional por membro da equipe. A Trinks anuncia R\$ 110 mensais para um a dois profissionais no ciclo mensal; a opção anual possui preço anunciado distinto. O Reservio permite 40 reservas em 30 dias no nível gratuito, com planos superiores para maior volume \cite{booksypricing,trinkspricing,reserviopricing}. Os valores são retratos das páginas consultadas e não cotações contratadas.

A diferenciação proposta pela Agendinha está na redução deliberada do escopo, no link público, na reserva sem conta para o cliente e na experiência instalável como PWA. Diferentemente dos concorrentes descritos, o projeto não inclui módulos de gestão financeira, estoque ou pagamentos. Não se afirma exclusividade dessas funcionalidades nem superioridade demonstrada: a efetiva vantagem em simplicidade dependerá da avaliação exploratória, e a econômica dependerá da apuração dos custos completos de operação.

\FloatBarrier
